% Standalone problem document, generated from the adjacent README.md.
% Compile with XeLaTeX. Mathematical content is maintained in Markdown.
\documentclass[11pt,a4paper]{article}
\usepackage[fontsize=10.5pt]{scrextend}
\usepackage[margin=22mm,headheight=15pt,headsep=9mm,footskip=11mm]{geometry}
\usepackage{fontspec}
\setmainfont{texgyrepagella}[Extension=.otf,UprightFont=*-regular,BoldFont=*-bold,ItalicFont=*-italic,BoldItalicFont=*-bolditalic]
\setsansfont{texgyreheros}[Extension=.otf,UprightFont=*-regular,BoldFont=*-bold,ItalicFont=*-italic,BoldItalicFont=*-bolditalic]
\setmonofont{texgyrecursor}[Extension=.otf,UprightFont=*-regular,BoldFont=*-bold,ItalicFont=*-italic,BoldItalicFont=*-bolditalic,Scale=MatchLowercase]
\usepackage{amsmath,amssymb}
\usepackage{unicode-math}
\setmathfont{texgyrepagella-math.otf}
\usepackage{xcolor,microtype,enumitem,needspace,fancyhdr}
\usepackage{bookmark}
\definecolor{ink}{HTML}{17344B}
\definecolor{accent}{HTML}{197D80}
\definecolor{muted}{HTML}{596773}
\hypersetup{colorlinks=true,linkcolor=accent,urlcolor=accent,pdftitle={IE-03 - Cryer's
Hadamard complete-pivoting
conjecture},pdfauthor={Open Problems in Numerical Linear Algebra}}
\urlstyle{same}
\setlength{\parindent}{0pt}
\setlength{\parskip}{5pt plus 1pt minus 1pt}
\linespread{1.04}
\setlength{\emergencystretch}{3em}
\widowpenalty=10000
\clubpenalty=10000
\displaywidowpenalty=10000
\allowdisplaybreaks[1]
\setlist{nosep,leftmargin=1.5em}
\providecommand{\tightlist}{\setlength{\itemsep}{0pt}\setlength{\parskip}{0pt}}
\providecommand{\pandocbounded}[1]{#1}
\pagestyle{fancy}
\fancyhf{}
\fancyhead[L]{\sffamily\footnotesize\color{muted}OPEN PROBLEMS IN NUMERICAL LINEAR ALGEBRA}
\fancyhead[R]{\sffamily\bfseries\footnotesize\color{accent}IE-03}
\fancyfoot[L]{\parbox[b]{0.9\textwidth}{\sffamily\fontsize{8}{10}\selectfont\color{muted}Editorial ratings; a bounded search cannot certify that a problem remains unsolved.\\Literature check: 2026-09-08}}
\fancyfoot[R]{\sffamily\footnotesize\color{muted}\thepage}
\renewcommand{\headrulewidth}{0.3pt}
\renewcommand{\headrule}{\hbox to\headwidth{\color{accent}\leaders\hrule height \headrulewidth\hfill}}
\makeatletter
\renewcommand{\section}{\Needspace{5\baselineskip}\@startsection{section}{1}{0pt}{12pt plus 2pt minus 2pt}{4pt}{\sffamily\bfseries\large\color{ink}}}
\renewcommand{\subsection}{\Needspace{4\baselineskip}\@startsection{subsection}{2}{0pt}{9pt plus 1pt minus 1pt}{3pt}{\sffamily\bfseries\normalsize\color{ink}}}
\makeatother
\setcounter{secnumdepth}{0}
\begin{document}
{\sffamily\small\color{muted}Linear systems and elimination\par}
\vspace{3pt}
{\raggedright\hyphenpenalty=10000\exhyphenpenalty=10000\sffamily\bfseries\fontsize{21}{24}\selectfont\color{ink}Cryer's
Hadamard complete-pivoting conjecture\par}
\vspace{7pt}
{\sffamily\small\textbf{Difficulty:} extreme\quad\textbf{Importance:} interesting
to the community\par}
\vspace{4pt}
{\color{accent}\hrule height 0.8pt}
\vspace{6pt}
\textbf{Status:} open.

\subsection{Context and notation}\label{context-and-notation}

All elimination in this problem is in exact arithmetic. Write
\(\|A\|_{\max}=\max_{ij}|a_{ij}|\). A pivoting path creates successive
active Schur complements \(S_1=A,S_2,\ldots,S_n\), with row/column
permutations as appropriate. Its element-growth factor is

\[
\rho(A)=\frac{\max_{1\leq j\leq n}\|S_j\|_{\max}}{\|A\|_{\max}}.
\]

Partial pivoting chooses a largest-magnitude entry in the active first
column; complete pivoting chooses one anywhere in the active matrix. If
ties occur, a universal statement includes every admissible tie choice;
a supremum includes all admissible paths. These conventions remove
implementation-dependent ambiguity. Matrices are nonsingular unless
otherwise stated.

\subsection{Problem statement}\label{problem-statement}

A real Hadamard matrix of order \(n\) is a matrix
\(H\in\{-1,1\}^{n\times n}\) satisfying \(HH^T=nI\). Prove or disprove
that, for every Hadamard matrix and every complete-pivoting path,

\[
\rho(H)=n.
\]

The quantified input is an existing Hadamard matrix; this question does
not ask whether Hadamard matrices exist in every order divisible by
four. Since \(\|H\|_{\max}=1\), the requested upper bound says that no
entry in any active Schur complement can exceed \(n\) in absolute value.
The final pivot already supplies the matching lower bound. Permuting or
resigning rows and columns does not remove the need to account for all
allowed pivot paths.

\subsection{References}\label{references}

N. J. Higham,
\href{https://doi.org/10.1137/1.9780898718027}{\emph{Accuracy and
Stability of Numerical Algorithms}, 2nd ed.}, SIAM (2002), Problem 9.17,
p.~193, poses this book-sourced question. Kravvaritis and Mitrouli,
\href{https://doi.org/10.1002/nla.637}{\emph{The growth factor of a
Hadamard matrix of order 16 is 16}}, NLA with Applications 16 (2009),
715--743, introduction and main result. Peca-Medlin,
\href{https://doi.org/10.1080/03081087.2026.2660796}{\emph{Complete
pivoting growth of butterfly matrices and butterfly Hadamard matrices}},
2026, introduction and §2.

\subsection{Status check}\label{status-check}

Searches for \textit{Hadamard Cryer 2026} and
\textit{Hadamard growth conjecture 2026 proof} found the general
conjecture still identified as open in the 2026 butterfly paper.
Shah--Urschel's August 2026 general elimination construction is not a
Hadamard counterexample. Known Sylvester and small-order cases must not
be counted as additional open problems.
\end{document}
